\documentclass[10pt,a4paper]{amsart}
\usepackage[margin=19mm]{geometry}
\usepackage{amsmath,amssymb,mathtools}
\usepackage[scr=rsfs,cal=euler]{mathalfa}
\usepackage{xfrac}
\usepackage[unicode,hidelinks,bookmarks=false]{hyperref}
\newcommand{\ie}{i.e.\ }
\newcommand{\cf}{cf.\ }
\newcommand{\resp}{respectively\ }
\newcommand{\Subsection}{\subsection}
\providecommand{\nobreakdash}{\nobreak\hbox{-}}
\setlength{\emergencystretch}{2em}
\multlinegap0pt
\usepackage{bm}
\usepackage{bbm}
\usepackage{dsfont}
\usepackage{xspace}
\usepackage{cancel}
\usepackage{mathtools}
\usepackage{amsthm}
\makeatletter
\@ifundefined{thm}{\newtheorem{thm}{Theorem}}{}
\makeatother
\def\cO{\mathcal{O}}
\def\cS{{\mathcal S}}
\def\div{\textrm{div}}
\title[Delta invariant of $\mathbb{Q}$-Cartier curve germs and the ]{Delta invariant of $\mathbb{Q}$-Cartier curve germs and the genus of representable numerical semigroups}
\author{Zsolt Baja, Tam\'as L\'aszl\'o, Andr\'as N\'emethi}
\date{Source: arXiv:2511.03406v1 (CC BY 4.0)}
\begin{document}
\maketitle

\noindent\emph{Source and licence.} Delta invariant of $\mathbb{Q}$-Cartier curve germs and the genus of representable numerical semigroups. Version arXiv:2511.03406v1 (CC BY 4.0), distributed under the Creative Commons Attribution 4.0 International licence, as recorded in the arXiv licence metadata for this exact version and shown on the abstract page https://arxiv.org/abs/2511.03406v1. Full rights notice: licenses/arxiv-2511.03406-CC-BY-4.0.txt.\par\smallskip

\noindent\textit{Editorial scope.} Counted unit: the Thm whose source label is \texttt{mainsmgps} in the article (source0.tex lines 797--799, source label \texttt{mainsmgps}), delivered together with the article's own complete original proof of it (source0.tex lines 802--830, one \texttt{proof} environment of 29 lines). No mathematics is written, repaired or re-derived: statement and proof are the article's own text line for line. Required context retained uncounted: none. Every definition, display and label of those lines is retained unchanged. Omitted from this package and not redistributed: 1--796, 800--801, 831--1155. Nothing inside a retained excerpt is omitted or altered: every retained body is delivered exactly as the article's lines are written, so no omission receipt of any kind accompanies this package. Citations: the retained excerpts carry no citation command at all, so no bibliography entry of the article is transcribed and no reference is redistributed. Reference adaptation: none was needed and none was made; every reference inside the delivered unit resolves inside it, so no unresolved reference or citation is produced. Printed slips kept literal: no printed word, symbol, hypothesis or number of the counted statement or of its proof was altered. Layout and class: the article's own class is \texttt{amsart} with packages including geometry, tikz, tkz-graph, float, pgfplots, inputenc, amsmath, amsthm, amsfonts, amssymb, mathtools, enumerate, graphicx, hyperref; the wrapper is the lane's pinned \texttt{amsart} editorial wrapper at 10pt on A4, which restates the theorem-like environments the retained text needs and adds the article's own macro definitions that the retained text needs (\texttt{\textbackslash cO}, \texttt{\textbackslash cS}, \texttt{\textbackslash div}), with extra wrapper packages (bm, bbm, dsfont, xspace, cancel, mathtools). The delivered numbering is the unit's own. Assets: no retained span contains a figure environment, a graphics-inclusion command, a PDF-page inclusion, a table or a TikZ picture; no image file is copied into this package, the package declares no asset dependency and no image file is redistributed with it. Licence: version arXiv:2511.03406v1 of this article is distributed under the Creative Commons Attribution 4.0 International licence, as recorded in the arXiv licence metadata for this exact version and shown on the abstract page https://arxiv.org/abs/2511.03406v1. A full rights notice is delivered at licenses/arxiv-2511.03406-CC-BY-4.0.txt. Characters: every retained excerpt is pure ASCII, so the delivered engine, XeLaTeX with its default Latin Modern text font, reports no missing character.
\begin{thm}\label{mainsmgps}
The numerical semigroup $\cS_{C^{gen}}$ associated with the germ of the irreducible curve $(C^{gen},0)$ coincides with $\cS_{\Gamma}$.   
\end{thm}

\begin{proof}
 We consider the canonical equivariant resolution $\pi$ as a good embedded resolution for $(C^{gen},0)\subset (X,0)$. We also write $p:=\tilde{C}^{gen}\cap E_0\in \tilde{X}$. 
 
 Recall that $\ell \in \cS_{C^{gen}}$ if and only if there is a function $f\in \cO_{C^{gen}}$ with $v(f)=\ell$, where $v$ is the valuation induced by the normalization of $(C^{gen},0)$. This condition is equivalent to the existence of a function $f\in \cO_{(X,0)}$ such that $(\div(f\circ \pi),\tilde C^{gen})_{\tilde{X}}=\ell$, where $\div(f\circ \pi):=(f)_{\Gamma}+\tilde C_f$ is the total transform,  $(f)_{\Gamma}$ is its part supported on $E=\pi^{-1}(0)$ and $\tilde C_f$ is the strict transform, and $(\ ,\, )_{\tilde{X}}$ denotes the divisorial intersection on $\tilde{X}$.
In particular, if the $E_0$ coefficient of $(f)_\Gamma$ is $m_{E_0}(f)$, then 
$\ell=v(f)=m_{E_0}(f)+(\tilde{C}_f, \tilde{C}^{gen})_p$, where the last term is the intersection multiplicity at $p\in\tilde{X}$. 

On the other hand, $\cS_{(X,0)}=\{\ell\in\mathbb{Z}_{\geq 0} | R_{X,\ell}\neq 0\}$ has the following reinterpretation. First note that if $f\in R_{X, \ell}$ is a homogeneous function, then 
$m_{E_0}(f)$ is exactly its homogeneous degree $\ell$. Hence  $\cS_{(X,0)}
= \{ m_{E_{0}}(f)\ |\ f \ \mbox{is homogeneous}\}$. 

The inclusion $\cS_{(X,0)}\subset  \cS_{C^{gen}}$ is rather direct. Let $f_{\ell}\in  R_{X,\ell}$
be a nonzero element of homogeneous degree $\ell$, and let $\tilde{C}_{f_{\ell}}$ be its strict transform. 
Then we can choose the generic orbit $\tilde{C}^{gen}$ in such a way that 
it does not intersect (along $E_0$) any of the strict transforms $\tilde{C}_{f_{\ell}}$, hence $(\tilde{C}_{f_{\ell}}, \tilde{C}^{gen})_p=0$.
Then $v$ associated with  $\tilde{C}^{gen}$ has value $\ell=m_{E_0}(f_\ell)=v(f_\ell)$.

Next, let us fix $\tilde{C}^{gen}$  and 
assume that $\ell\in  \cS_{C^{gen}}$. Choose  a function $f\in \cO_{(X,0)}$ such that 
$\ell=v(f)$. Write $f$ as a sum of  its homogeneous components $f=\sum _{j\geq 0} f_j$, where 
each $f_j \in R_{X,j}$ is nonzero. Then, for any fixed $j$, the strict transform $\tilde{C}_{f_j}$
of $f_j$ is a finite union of $\mathbb{C}^*$-orbits  intersected with $\tilde{X}$. 
It might contain the fixed orbit $\tilde{C}^{gen}$ or not. 
Therefore, $v(f_j)=j+(\tilde{C}_{f_j}, \tilde{C}^{gen})_p$ equals $j$ if 
$\tilde{C}_{f_j}$ does not contain the fixed orbit $\tilde{C}^{gen}$, and it is $\infty$ otherwise. 
In particular, all the possible finite values $\{v(f_j)\}_j$ are distinct, and 
$\ell=v(f)=\min \{v(f_j)\ |\  v(f_j)<\infty\}$. This shows that $\ell=v(f)=v(f_j)=j$
for some $f_j\in R_{(X,0)}$, hence $\ell\in \cS_{(X,0)} $. 
\end{proof}

\end{document}
